\documentclass[11pt]{article}
\usepackage{amsmath,amssymb,amsthm}
\usepackage[margin=1in]{geometry}
\usepackage{hyperref}
\usepackage{array}
\usepackage{booktabs}

\newtheorem{theorem}{Theorem}
\newtheorem{conjecture}[theorem]{Conjecture}
\newtheorem{proposition}[theorem]{Proposition}
\newtheorem{lemma}[theorem]{Lemma}
\newtheorem{corollary}[theorem]{Corollary}
\newtheorem{observation}[theorem]{Observation}
\theoremstyle{definition}
\newtheorem{definition}[theorem]{Definition}
\newtheorem{example}[theorem]{Example}
\theoremstyle{remark}
\newtheorem{remark}[theorem]{Remark}

\newcommand{\Aalt}{A^{\mathrm{alt}}}
\newcommand{\thetaalt}{\theta^{\mathrm{alt}}}
\newcommand{\wt}{\mathrm{wt}}
\newcommand{\Sn}{\mathfrak{S}_n}
\newcommand{\ZZ}{\mathbb{Z}}
\newcommand{\QQ}{\mathbb{Q}}
\newcommand{\NN}{\mathbb{N}}
\newcommand{\Bru}{\leq_B}
\newcommand{\qint}[1]{[#1]_t}

\title{The block-triangular structure of nil-Hecke atoms:\\
       six cases verified, with a refinement to the\\
       two-parameter Gaussian conjecture}
\author{Clio (Anthropic)}
\date{2026-07-25}

\begin{document}
\maketitle

\begin{abstract}
We extend the verification of Route~$\delta^+$
(cf.\ \texttt{2026-07-25-t-graded-A-G-verification-4-cases.pdf})
from four cases to six by adding $\mu = (3,2,1)$ at $n=3$ and
$\mu = (2,2,2,0)$ at $n=4$. The block-lower-triangular structure of
the nil-Hecke atoms $\Aalt_\gamma$ expressed over crystal elements
holds in both new cases, and the global identity
$\sum_\gamma c_\gamma(t)\, \Aalt_\gamma = P_\mu(x;t)$ is verified
symbolically. The two-parameter Gaussian conjecture
$c_\mu = \qint{n} \cdot [k]_{t^m}$ (with $m = |\mathrm{Stab}(\mu)|$)
that had been floated on the strength of $\mu = (2,2,0,0)$ is
\emph{falsified} by $\mu = (2,2,2,0)$: the stabiliser has order six,
yet no $t^m$-flavoured factor appears. We replace it with a cleaner
dichotomy indexed by the Bruhat-order geometry of the orbit:
when $\mathrm{orbit}(\mu)$ is a total Bruhat chain, $c_\gamma$
is the pure Gaussian $\qint{L(\gamma)+1}$; when Bruhat-length
collisions occur and their partners are related by the residual
action of the stabiliser, $[k]_{t^m}$-type corrections appear.
The mod-$t$ structure remains exact in all six cases, so the Dim
Lemma stays unconditional.
\end{abstract}

\section{Introduction}

Fix $n\geq 2$ and a composition $\mu \vdash n$ (thought of as a
weakly decreasing weight, or, in the orbit view, as any element of
$\Sn \cdot \mu$). The nil-Hecke atom operators
\[
   \thetaalt_i(f) \;=\; \frac{x_{i+1} - t\, x_i}{x_i - x_{i+1}}
                        \bigl(f - s_i(f)\bigr),
   \qquad 1 \leq i \leq n-1,
\]
are polynomial: the numerator $x_{i+1} - t\, x_i$ absorbs the
divisibility of $f - s_i(f)$ by $x_i - x_{i+1}$. At $t = 0$ they
recover $x_{i+1}\, \partial_i$ (Mason's classical atom operators);
at $t = 1$ they collapse the Hecke bracket to zero.

For $\gamma \in \Sn \cdot \mu$, pick any shortest word
$s_{i_1} \cdots s_{i_r}$ carrying $\mu$ to $\gamma$ in the weak
order and set
\[
   \Aalt_\gamma \;=\; \thetaalt_{i_1} \cdots \thetaalt_{i_r} (x^\mu).
\]
(The braid relations for $\thetaalt$ ensure the result is
independent of the reduced word chosen.) These are the
\emph{nil-Hecke atoms}. The starting fact behind the sprint is
the atom decomposition
\begin{equation}\label{eq:atom-decomp}
   P_\mu(x;t) \;=\; \sum_{\gamma \in \mathrm{orbit}(\mu)}
                    c_\gamma(t)\, \Aalt_\gamma,
   \qquad c_\gamma(t) \in \ZZ_{\geq 0}[t],
\end{equation}
proven algebraically on 2026-07-23 conditional on the $V_{<\mu}$
consistency lemma and closed unconditionally on 2026-07-25 via
Route~$\delta^+$'s block-triangular mechanism.

Route~$\delta^+$ conjectures that, expressed over the highest-weight
crystal $B(\mu)$ (equivalently the standard basis of the polynomial
representation supported on shapes $\leq \mu$), the matrix
$\bigl[\Aalt_\gamma\bigr]_{\gamma,\, b}$ is
\begin{itemize}
\item block-lower-triangular in the Bruhat order on
      $\mathrm{orbit}(\mu)$, with rows indexed by $\gamma$ and columns
      by crystal elements grouped by classical atom;
\item on the diagonal block: $1$ on extremal elements,
      $(1-t)$ on the remaining classical-atom elements;
\item on strictly sub-diagonal blocks: pure-$t$ ``shadows'';
\item on strictly above-diagonal positions: entries divisible by $t$
      (the refinement forced on 2026-07-25 by $\mu = (2,2,0,0)$).
\end{itemize}

The purpose of this note is to add two data points --- one with a
completely distinct-parts orbit at $n=3$, one with a strongly
constrained orbit at $n=4$ --- and to draw the natural conclusion
about the shape of $c_\gamma(t)$.

\section{The two new cases}

Before turning to the cases proper, we recall the computational
setup used in both this note and the 4-cases predecessor. The
crystal $B(\mu)$ is realised inside the polynomial ring
$\ZZ[t][x_1, \ldots, x_n]$ as the span of monomials
$x^\nu$ with $\nu$ ranging over multisets that can appear on
the branches of the Hall--Littlewood $P_\mu$. In practice, for a
computer-algebra verification, one enumerates the semistandard
Young tableaux of shape $\mu$ (equivalently, the elements of
$B(\mu)$) and records their weight compositions $\nu$. The atoms
$\Aalt_\gamma$ are then computed by applying $\thetaalt_i$ symbolically
starting from $x^\mu$, using a stored reduced word to $\gamma$.
The block-triangular matrix is the coefficient array of
$\{\Aalt_\gamma\}$ over the standard monomial basis, sorted by
crystal element (grouped by classical atom $A_\gamma$).

We recall the diagnostic vocabulary from the 4-cases note.
\emph{Weak form}: for each $\gamma$, the support of $\Aalt_\gamma$
(over crystal elements) lies within the Bruhat down-set of $\gamma$
together with the diagonal block at $\gamma$. \emph{Strong form}:
the diagonal block equals $(1, 1-t, \ldots, 1-t)$ on the classical
atom $A_\gamma$; the sub-diagonal shadows lie in $t\, \ZZ[t]$;
the above-diagonal entries lie in $t\, \ZZ[t]$.

\subsection{Case A: $\mu = (3,2,1)$, $n = 3$}

All three parts distinct: $|\mathrm{orbit}(\mu)| = 6$, no stabiliser.
The Bruhat poset on the orbit is the full weak order on $\Sn$ at
$n=3$, with lengths $0,1,1,2,2,3$.

\begin{proposition}\label{prop:321}
For $\mu = (3,2,1)$:
\begin{enumerate}
\item[(i)]  All $6$ weak-form checks pass; all $6$ diagonal-block
            entries pass; all $13$ sub-diagonal shadow checks pass.
\item[(ii)] All above-diagonal entries are divisible by $t$.
            The extremal row $\gamma = (1,2,3)$ has non-monomial
            above-diagonal shadows: entry $t(t-1)$ at the weight
            $(2,2,2)$ column and $-t^2(t+1)$ at the $(3,2,1)$ column
            (both $\in t\,\ZZ[t]$).
\item[(iii)]The global identity
            $\sum_\gamma c_\gamma(t)\,\Aalt_\gamma = P_\mu(x;t)$
            holds symbolically in $\QQ(t)[x_1, x_2, x_3]$.
\end{enumerate}
\end{proposition}

The coefficient values assemble into
\begin{equation}\label{eq:c-321}
   \begin{aligned}
   L=0:\quad & c_{(3,2,1)} = \qint{2}\,\qint{3}, \\
   L=1:\quad & c_{(2,3,1)} = \qint{3}, \qquad c_{(3,1,2)} = \qint{2}^2, \\
   L=2:\quad & c_{(1,3,2)} = c_{(2,1,3)} = \qint{2}, \\
   L=3:\quad & c_{(1,2,3)} = 1.
   \end{aligned}
\end{equation}

\begin{remark}[Asymmetry within Bruhat length one]
The two Bruhat-length-$1$ orbit elements $(2,3,1)$ and $(3,1,2)$
are \emph{not} conjugate under any obvious symmetry available to
the coefficient function, and they receive genuinely different
Gaussians: $\qint{3}$ vs.\ $\qint{2}^2$. This is qualitatively the
same phenomenon Clio flagged on 07-25 for $\mu = (2,1,0)$
(different length-$1$ elements $(2,0,1)$ and $(1,2,0)$ received
$\qint{2}^2$ and $\qint{3}$ respectively). Bruhat length alone
does not determine $c_\gamma$: some position-dependent piece is at
work, and it lives entirely at length $1$ in these two examples.
The higher lengths ($L=2$: both $\qint{2}$; $L=3$: $1$) are
symmetric.
\end{remark}

To make the asymmetry concrete: the two Bruhat-length-$1$ orbit
elements $(2,3,1)$ and $(3,1,2)$ arise from $(3,2,1)$ by applying
$s_2$ and $s_1$ respectively. In each case the resulting nil-Hecke
atom, expressed over the crystal basis, occupies a $2$-element
classical atom (extremal $+$ one interior weight). The two
diagonal blocks have the same shape $(1, 1-t)$; the two
sub-diagonal shadow patterns \emph{differ}, because the reduced
words that generate the $2$-element classical atoms encounter
different collections of intermediate weights inside $B(\mu)$.
Symbolically,
\[
   \Aalt_{(2,3,1)} \;=\; \thetaalt_2(x^{(3,2,1)}), \qquad
   \Aalt_{(3,1,2)} \;=\; \thetaalt_1(x^{(3,2,1)}),
\]
and the two operators have different denominators
($x_2 - x_3$ vs.\ $x_1 - x_2$) that produce different
weight-decompositions. The upshot for $c_\gamma$ is exactly the
asymmetry $\qint{3}$ vs.\ $\qint{2}^2$: the same-degree Gaussians
$1 + t + t^2$ and $(1+t)^2 = 1 + 2t + t^2$ differ only in the
middle coefficient, and that middle coefficient is where the
position-dependence of the shadow pattern is recorded.

\begin{remark}[On the extremal shadow]
The extremal row --- the row corresponding to $\gamma = (1,2,3)$,
the Bruhat-maximal element --- deserves particular note. Under the
weak form, one might expect the extremal row to be a coordinate
vector; under the strong form, one might expect its off-diagonal
entries to lie in $t\, \ZZ_{\geq 0}[t]$. The value $-t^2(t+1)$ at
weight $(3,2,1)$ shows that the signs are not fixed: above-diagonal
entries are $t$-divisible, but not necessarily $t$-monomials with
non-negative coefficients. This is consistent with the 4-cases note
observation at $(2,2,0,0)$: sign cancellations in the sum
$\sum c_\gamma \Aalt_\gamma$ recover the (positive) monomial
symmetric expansion $\sum m_\lambda K_{\lambda\mu}(t)$.
\end{remark}

\subsection{Case B: $\mu = (2,2,2,0)$, $n = 4$}

Stabiliser $\mathfrak{S}_3 \times \mathfrak{S}_1$ (order $6$),
orbit size $4$. The four orbit elements are
$(2,2,2,0), (2,2,0,2), (2,0,2,2), (0,2,2,2)$; the orbit forms a
\emph{total} chain in Bruhat order, with lengths $0,1,2,3$.

\begin{proposition}\label{prop:2220}
For $\mu = (2,2,2,0)$:
\begin{enumerate}
\item[(i)]  All $4$ weak-form checks pass; all $10$ diagonal-block
            entries pass; all $10$ sub-diagonal shadow checks pass.
\item[(ii)] The global identity
            $\sum_\gamma c_\gamma(t)\,\Aalt_\gamma = P_\mu(x;t)$
            holds symbolically.
\end{enumerate}
\end{proposition}

The coefficient values form a pure Gaussian chain:
\begin{equation}\label{eq:c-2220}
   c_{(2,2,2,0)} = \qint{4}, \qquad
   c_{(2,2,0,2)} = \qint{3}, \qquad
   c_{(2,0,2,2)} = \qint{2}, \qquad
   c_{(0,2,2,2)} = 1.
\end{equation}

\begin{observation}\label{obs:no-tm}
The stabiliser of $\mu = (2,2,2,0)$ has order $6$, and yet no
factor of $[k]_{t^m}$ for any $m$ appears in $c_\gamma$. The four
coefficients are Gaussians in $t$ (not in any $t^m$ with $m > 1$),
one for each Bruhat length. Moreover $c_\gamma = \qint{L(\gamma)+1}$
throughout, reproducing the shape already seen for
$\mu = (2,2,0)$ and $\mu = (2,2,1)$ at $n = 3$.
\end{observation}

\begin{example}[The explicit chain identity at $(2,2,2,0)$]
Substituting \eqref{eq:c-2220} into \eqref{eq:atom-decomp} we get
\[
   P_{(2,2,2,0)}(x_1,\ldots,x_4;t) \;=\;
     \qint{4}\, \Aalt_{(2,2,2,0)}
     + \qint{3}\, \Aalt_{(2,2,0,2)}
     + \qint{2}\, \Aalt_{(2,0,2,2)}
     + \Aalt_{(0,2,2,2)}.
\]
Each atom on the right is easily computed: the extremal
$\Aalt_{(0,2,2,2)}$ is the monomial $x_2^2 x_3^2 x_4^2$; the
next $\Aalt_{(2,0,2,2)}$ is $\thetaalt_1$ applied to
$x_1^2 x_2^2 x_3^2$, which unfolds to
$x_1^2 x_3^2 x_4^2 + (1-t)(\text{shorter atom elements})$; and so
on up the chain. The coefficient chain
$1, \qint{2}, \qint{3}, \qint{4}$ is exactly what the sub-diagonal
$-t$-shadow recursion of the block-triangular matrix predicts,
starting from $c_{\text{top}} = 1$ and propagating downward.
\end{example}

\begin{remark}[Why $(2,2,2,0)$ is a Bruhat chain]
The reason $\mathrm{orbit}\bigl((2,2,2,0)\bigr)$ collapses to a
chain, despite the ambient Weyl group $\mathfrak{S}_4$ being
non-abelian and of order $24$, is the stabiliser. The four
orbit elements $(2,2,2,0), (2,2,0,2), (2,0,2,2), (0,2,2,2)$
can be encoded by the position of the single $0$: at slot $4, 3,
2, 1$ respectively. In this encoding, the Bruhat order on the
orbit is the transitive closure of the covering relation
``move the $0$ left by one slot'', which is a total chain. There
is no room in a single-zero configuration for two orbit elements
to have the same Bruhat length.
\end{remark}

\subsection{Comparative summary of the two new cases}

The two new cases occupy opposite ends of the spectrum available
at this size. $(3,2,1)$ maximises orbit size at $n = 3$
($|\mathrm{orbit}| = 6$, the full symmetric group $\Sn$), with
no stabiliser and hence with the largest possible Bruhat
combinatorics on the orbit; it also exhibits the most vivid
form of position-dependent asymmetry within a single Bruhat
length. $(2,2,2,0)$ minimises orbit size at $n = 4$ subject to
being non-trivial ($|\mathrm{orbit}| = 4$), because the large
stabiliser $\mathfrak{S}_3 \times \mathfrak{S}_1$ collapses most
of the ambient $\mathfrak{S}_4$; the resulting orbit is forced
into a single Bruhat chain, and $c_\gamma$ is forced into a
pure Gaussian sequence with no asymmetry.

The two cases together demonstrate that neither Bruhat length
alone (Case A shows this within a single length) nor stabiliser
size alone (Case B shows a large stabiliser producing the
cleanest structure) is the operative variable for $c_\gamma$.
What matters is the geometry of the orbit as a Bruhat poset, and
in particular whether Bruhat lengths collide there.

\section{Retirement of the two-parameter Gaussian; a Bruhat-chain
         replacement}

The 4-cases note ended by floating, on the basis of
$\mu = (2,2,0,0)$, the tentative conjecture
\begin{equation}\label{eq:old-conj}
   c_\mu \;\stackrel{?}{=}\; \qint{n} \cdot [k]_{t^m},
   \qquad m = |\mathrm{Stab}(\mu)|.
\end{equation}
For $(2,2,0,0)$, $n=4$, $\mathrm{Stab}$ of order $2$, we saw
$c_{(2,2,0,0)} = \qint{3} \cdot [2]_{t^2}$, and this
\emph{looked like} evidence for \eqref{eq:old-conj}. Observation
\ref{obs:no-tm} falsifies \eqref{eq:old-conj}: for
$(2,2,2,0)$, $n=4$, $|\mathrm{Stab}| = 6$, we get $\qint{4}$,
not $\qint{4} \cdot [k]_{t^6}$.

The corrected picture separates two regimes according to the
Bruhat geometry of $\mathrm{orbit}(\mu)$:

\begin{conjecture}[Bruhat-chain / Bruhat-collision dichotomy]
\label{conj:dichotomy}
Let $\mu$ be a composition with orbit $\mathrm{orbit}(\mu) \subset
\Sn \cdot \mu$ carrying the Bruhat order inherited from $\Sn$.
\begin{enumerate}
\item[(a)] \emph{Chain regime.} If $\mathrm{orbit}(\mu)$ is a total
           Bruhat chain --- i.e.\ each length
           $0, 1, \ldots, |\mathrm{orbit}(\mu)| - 1$ occurs
           exactly once --- then
           \[
              c_\gamma(t) \;=\; \qint{L(\gamma) + 1},
           \]
           the pure Gaussian in $t$ of degree $L(\gamma)$.
\item[(b)] \emph{Collision regime.} If Bruhat lengths collide on
           $\mathrm{orbit}(\mu)$ and the colliding partners are
           related by the residual action of $\mathrm{Stab}(\mu)$,
           then $c_\gamma$ carries additional $[k]_{t^m}$-type
           factors keyed to the length of the residual orbit.
\end{enumerate}
\end{conjecture}

Case B, $(2,2,2,0)$, is a clean instance of (a): the orbit is a
$4$-element chain of lengths $0, 1, 2, 3$ and \eqref{eq:c-2220}
is exactly $\qint{L+1}$. Cases $\mu = (2,2,0)$, $(2,2,1)$ of the
4-cases note are also (a) (they are chains of lengths $0,1$ and
$0, 1, 2$).

The case $\mu = (2,2,0,0)$ is (b): both $(2,0,0,2)$ and
$(0,2,2,0)$ have Bruhat length $2$ and are exchanged by an element
of the residual action, producing the $[2]_{t^2}$ factor. That
residual action lives in $\mathfrak{S}_4 / (\mathrm{Stab})$, and
the two length-$2$ orbit elements are paired inside a single
residual class, which is the mechanism the conjecture proposes as
the source of the $[2]_{t^m}$ factor with $m = 2$
(equal to the length of the residual orbit).

Cases $\mu = (2,1,0)$ and $\mu = (3,2,1)$ are simultaneous
collision-and-asymmetry cases: lengths collide (two elements of
length $1$; two elements of length $2$) but with a trivial
stabiliser, so the ``collision partners'' are not literally
identified by the stabiliser. Here Conjecture~\ref{conj:dichotomy}
does not immediately predict the shape; the observed values
\eqref{eq:c-321} include the anomalous $\qint{2}^2 = (1+t)^2$
term at length $1$, position $(3,1,2)$. Understanding this
position-dependent asymmetry --- and where inside the collision
regime it lives --- is the natural next probe.

\begin{remark}[Testable predictions]
The dichotomy is falsifiable at low $n$. The compositions with
Bruhat-collision orbits and non-trivial stabilisers that Clio has
\emph{not} yet computed include $(3,3,1,0)$, $(3,3,2,0)$,
$(3,2,2,0)$, and $(2,2,1,1)$. If Conjecture~\ref{conj:dichotomy}
holds, the first three should exhibit $[k]_{t^m}$ factors keyed
to the length-$2$ residual orbit; $(2,2,1,1)$ has the largest
stabiliser and should be the most stringent test.
\end{remark}

\section{Consequences and the status of Route~$\delta^+$}

The six cases now on hand are summarised in Table~\ref{tab:six}.

\begin{table}[h]
\centering
\renewcommand{\arraystretch}{1.3}
\begin{tabular}{lcccl}
\toprule
$\mu$ & $|\mathrm{orbit}|$ & Stabiliser & Chain? &
     $c_\gamma$ pattern \\
\midrule
$(2,2,0)$    & $3$ & $\langle s_1\rangle$ & yes &
     $1, \qint{2}, \qint{3}$ (Gaussian chain) \\
$(2,2,1)$    & $3$ & $\langle s_1\rangle$ & yes &
     $1, \qint{2}, \qint{3}$ (Gaussian chain) \\
$(2,1,0)$    & $6$ & trivial & no &
     collision + asymmetry at $L=1$ \\
$(3,2,1)$    & $6$ & trivial & no &
     collision + asymmetry at $L=1$ \\
$(2,2,0,0)$  & $6$ & $\langle s_1\rangle$ & no &
     collision with $[2]_{t^2}$ (case b) \\
$(2,2,2,0)$  & $4$ & $\mathfrak{S}_3 \times \mathfrak{S}_1$ & yes &
     $1, \qint{2}, \qint{3}, \qint{4}$ (chain) \\
\bottomrule
\end{tabular}
\caption{The six cases now verified. ``Chain?'' asks whether
$\mathrm{orbit}(\mu)$ is a total chain in Bruhat order; when yes,
$c_\gamma = \qint{L(\gamma)+1}$.}
\label{tab:six}
\end{table}

Three structural conclusions:

\emph{The Bruhat-chain regime is a subordinate tautology.} Where
the orbit is a total Bruhat chain, $c_\gamma = \qint{L+1}$ is
essentially forced: the block-triangular structure plus the
$(1-t)$ diagonal plus the top-of-chain boundary condition
$c_{\text{top}} = 1$ propagates up the chain by a single first-order
recursion in $t$. Any structural proof of Route~$\delta^+$ that
does not \emph{start} by handling the chain case is doing the wrong
thing; conversely, the chain case is the target for the simplest
first proof.

To make the chain recursion explicit: let
$\gamma_0 <_B \gamma_1 <_B \cdots <_B \gamma_N$ be the chain, with
$L(\gamma_k) = k$, and let $c_k \in \ZZ[t]$ denote $c_{\gamma_k}$.
Row $\gamma_k$ of the block-triangular matrix (over the crystal
elements) has diagonal entry $1 - t$ (or $1$ at $k = 0$) and a
single sub-diagonal $t$-shadow into $\gamma_{k-1}$ of value $-t$
(the only shadow possible in a chain). The requirement that the
row-combination $\sum_k c_k \Aalt_{\gamma_k}$ evaluate to a
symmetric polynomial forces
\[
   c_k \cdot (1 - t) + c_{k-1} \cdot (-t) \cdot [\text{shadow sign}]
     \;=\; \text{diagonal entry of the symmetric polynomial},
\]
which under the top-of-chain boundary $c_0 = 1$ yields
$c_k = 1 + t + t^2 + \cdots + t^k = \qint{k+1}$. This is the
recursion visible in the chain data of Table~\ref{tab:six}.

\emph{Collision regime is where the content is.} Both categorical
frameworks (the crystal-side quasicrystal contraction of
Brauner--Daugherty--Mason--Schilling and the singular Soergel
bimodule polynomial-forcing of Elias--Ko--Libedinsky--Patimo, see
below) must specialise to \emph{some} statement about how residual
symmetry organises length-colliding orbit elements. The
$[2]_{t^2}$ appearance for $(2,2,0,0)$ is the first instance of
that; the second instance --- where the residual acts by a group
of order $3$ or higher --- is the diagnostic test for
$(3,3,1,0)$-type cases. The trivial-stabiliser
collision-and-asymmetry cases $(2,1,0)$ and $(3,2,1)$ are a
third, distinct sub-regime in which the collision persists but
no group-theoretic mechanism identifies the partners; that
sub-regime is the least well-understood at present, and its
appearance in the smallest possible non-chain examples is one of
the surprises of the six-case data.

\emph{Mod-$t$ is exact in all six cases.} Setting $t = 0$ in the
diagonal block collapses $(1-t) \mapsto 1$; the sub-diagonal and
above-diagonal $t$-shadows vanish. The classical Mason atom
decomposition $B(\mu) = \bigsqcup_\gamma A_\gamma$
(Lascoux--Sch\"utzenberger 1990, Mason 2006) is recovered, and the
Dim Lemma continues to require only this classical structure ---
it remains unconditional. In particular, the observation that the
$c_\gamma$ evaluated at $t = 0$ all equal $1$ (both in the chain
regime $\qint{L+1}|_{t=0} = 1$ and in the collision regime
$\qint{n} \cdot [k]_{t^m}|_{t=0} = 1$) is not accidental: it is
the statement that at the classical limit the Schur function
$s_\mu$ decomposes as the sum of classical Demazure atoms indexed
by the orbit, each with multiplicity one. The $t$-graded content
of Route~$\delta^+$ lives entirely in the positive powers of $t$,
where the various Gaussians distinguish themselves.

\section{Categorical shadows}

Two categorifications, both landed in the 2026-07-24-bis browse
cycle, are candidate homes for a structural proof of
Conjecture~\ref{conj:dichotomy}.

\emph{The crystal side: Brauner--Daugherty--Mason--Schilling
2607.12232.} This paper gives a rigorous \emph{classical} ($t=0$)
Bruhat contraction of the crystal skeleton on unions of Demazure
crystals. Their tiles are the Young quasisymmetric Schur functions
$\mathrm{YQS}_\alpha$, which are coarser than the Mason atoms
$A_\gamma$ that appear here (they collect multiple atoms into a
single quasisymmetric block indexed by a composition rearrangement).
The $t = 0$ specialisation of Route~$\delta^+$'s block-triangular
matrix should therefore refine their tiled contraction. What is
missing from their picture --- and what any structural proof of
$\delta^+$ must supply --- is the $t$-graded lift: the
$[k]_{t^m}$ factors that Conjecture~\ref{conj:dichotomy}(b)
predicts do not appear anywhere in a classical crystal skeleton
and must come from a $t$-graded refinement of the contraction.

\emph{The Soergel side: Elias--Ko--Libedinsky--Patimo 2407.13128.}
The Atomic Leibniz Rule is a polynomial-forcing statement for
singular Soergel bimodules indexed by double-coset combinatorics
$W_I \backslash W / W_J$. The name collision with the sprint's
nil-Hecke atoms is not (yet) a proved equivalence, but the shape of
the singular double-coset table is precisely the residual-action
combinatorics that Conjecture~\ref{conj:dichotomy}(b) invokes.
Reading 2407.13128 for whether its polynomial-forcing table produces
$[k]_{t^m}$ factors of the shape observed here is a deep-read
priority.

Neither framework has been read yet at the level of detail needed
to convert either observation into a proof.

\emph{The Fock side (Robin's picture).} The transfer-operator
computations that Robin has been developing on the Fock side, in
which $c^\lambda_{\mu\nu}$ arises through Kostka-matrix inversion,
provide the fourth face of the polytope. On that side, the
determinism of the Kostka inversion is the analogue of the
uniqueness statement in the Dim Lemma; the Bruhat-chain vs.\
Bruhat-collision dichotomy should have a Fock-side avatar in the
shape of the Kostka matrix relative to the residual action of
$\mathrm{Stab}(\mu)$. Whether the two dichotomies are literally
the same statement in different notation, or are two features of
a single underlying combinatorial object, is one of the sharper
open questions left after these six cases.

The four faces of the conjectural polytope --- Fock (Robin),
atoms (Clio), crystals (Brauner--Daugherty--Mason--Schilling),
Soergel bimodules (Elias--Ko--Libedinsky--Patimo) --- each
propose a different vocabulary for the same block-triangular
structure. What Route~$\delta^+$'s six-case data isolate is a
minimum common statement that any of the four vocabularies must
be able to prove: the chain-and-collision dichotomy on
$\mathrm{orbit}(\mu)$, with the specific coefficient values
recorded in \eqref{eq:c-321} and \eqref{eq:c-2220} above. If any
one of the four frameworks can produce those coefficient values
from its native combinatorics, Route~$\delta^+$ has found its
naming home.

\section{Next steps}

\begin{enumerate}
\item \emph{Bruhat-collision probes with larger stabilisers.}
      Compute $c_\gamma$ for $(3,3,1,0)$, $(3,3,2,0)$, $(3,2,2,0)$,
      and $(2,2,1,1)$ and check whether $[k]_{t^m}$ factors appear
      with $m$ predicted by the residual-orbit length. The
      $(2,2,1,1)$ case, with stabiliser
      $\mathfrak{S}_2 \times \mathfrak{S}_2$, is the most stringent
      test.
\item \emph{Structural proof of the Bruhat-chain case.} Attempt a
      direct proof of Conjecture~\ref{conj:dichotomy}(a) from the
      block-triangular structure alone. The Gaussian recursion
      $\qint{L+1} = 1 + t\,\qint{L}$ should match the sub-diagonal
      $t$-shadow contribution to the recursion in $L$.
\item \emph{Deep-read Elias--Ko--Libedinsky--Patimo 2407.13128}
      for singular double-coset polynomial forcing, checking
      whether the table produces the $[k]_{t^m}$ pattern seen at
      $(2,2,0,0)$.
\end{enumerate}

\begin{remark}[On why $(2,2,1,1)$ is the sharpest test]
Compositions $(2,2,1,1)$ has stabiliser $\mathfrak{S}_2 \times
\mathfrak{S}_2$ of order $4$; the orbit has size $6$; Bruhat
lengths on the orbit include a collision at length~$1$ (the two
adjacent swaps of a $2$-block against a $1$-block). If the
dichotomy holds and the residual action pairs those two length-$1$
elements into a single residual class of size $2$, we should see a
$[2]_{t^2}$-type factor at $c_{(2,2,1,1)}$. If instead the
stabiliser factor gives a $[k]_{t^m}$ with $m > 2$ (say $m = 4$
coming from the full stabiliser order), the residual-orbit-length
formulation is wrong and only the raw stabiliser size matters.
Either outcome is informative; the current data do not
distinguish them.
\end{remark}

\section*{What the six-case pattern reveals}

The chain / collision dichotomy is not something we would have
guessed from $(2,2,0)$ alone. It only becomes visible once one has
in hand \emph{both} a stabiliser-rich chain (like $(2,2,2,0)$) and
a stabiliser-poor collision (like $(2,1,0)$ or $(3,2,1)$): each
isolates one of the two features that the earlier two-parameter
Gaussian conjecture had accidentally conflated. What appears to
matter is not the size of the stabiliser but whether Bruhat lengths
collide on the orbit at all, and if so whether the residual
stabiliser action identifies the collision partners. The object
under study --- $c_\gamma(t)$ --- is beginning to look like a
genuinely two-story invariant: a Gaussian propagation along Bruhat
chains, decorated by a residual-symmetry correction wherever the
Bruhat order permits.

\end{document}
